\documentclass[11pt]{article}

\usepackage[utf8]{inputenc}
\usepackage[T1]{fontenc}
\usepackage[a4paper, margin=0.9in]{geometry}
\usepackage{booktabs}
\usepackage{longtable}
\usepackage{array}
\usepackage{enumitem}
\usepackage{xcolor}
\usepackage{fancyhdr}
\usepackage{hyperref}

\definecolor{ok}{RGB}{16,122,86}
\definecolor{warn}{RGB}{176,110,10}
\definecolor{slate}{RGB}{70,80,100}

\hypersetup{colorlinks=true, linkcolor=ok, urlcolor=ok}

\pagestyle{fancy}
\fancyhf{}
\lhead{BookVault}
\rhead{Demonstration \& Viva Guide}
\cfoot{\thepage}

\setlength{\parindent}{0pt}
\setlength{\parskip}{5pt}

\newcommand{\proof}[1]{\textcolor{ok}{\textbf{Proves:}} #1}
\newcommand{\note}[1]{\textcolor{warn}{\textbf{Note:}} #1}

\title{\textbf{BookVault} \\[4pt]
\large Demonstration \& Viva Guide}
\author{CSE447 Lab Project $\mid$ Choice 2}
\date{}

\begin{document}

\maketitle
\thispagestyle{fancy}

\begin{center}
\fbox{\begin{minipage}{0.92\textwidth}
\textbf{How to use this guide.} Part A and B are click-by-click walkthroughs
(patron, then librarian). Part C shows how to demonstrate each of the nine
security features. Part D maps every lab requirement to where it lives.
Part E proves the cryptography from the terminal. Part F is a list of likely
viva questions with short answers.
\end{minipage}}
\end{center}

% ==========================================================
\section{Before you start}

Open a terminal in the project folder and run:

\begin{verbatim}
   npm run dev
\end{verbatim}

This starts both parts of the system:

\begin{itemize}[leftmargin=1.4cm, itemsep=2pt]
  \item \textbf{API} (Express backend) on \texttt{http://localhost:4000}
  \item \textbf{App} (React frontend) on \texttt{http://localhost:5173}
\end{itemize}

Open \texttt{http://localhost:5173} in a browser.

\subsection*{Accounts}

\begin{center}
\begin{tabular}{lll}
\toprule
\textbf{Role} & \textbf{Username} & \textbf{Password} \\
\midrule
Head Librarian & \texttt{librarian} & \texttt{Vault-Admin-2026} \\
Patron & \texttt{faiza} & \texttt{Reading-Room-77} \\
\bottomrule
\end{tabular}
\end{center}

\subsection*{Signing in takes two steps}

\begin{enumerate}[leftmargin=1.4cm, itemsep=2pt]
  \item Enter username and password, press \textbf{Continue}.
  \item A six-digit code is asked for. Click \textbf{Show current code},
        then type those six digits.
\end{enumerate}

\note{Only one session per account is live at a time. Signing in as the
librarian signs the patron out. To show both at once, use two different
browsers (for example Chrome and Firefox), not two tabs.}

% ==========================================================
\section{Part A --- Patron walkthrough}

Sign in as \texttt{faiza}.

\subsection*{F1. Registration \& Login}

Sign out and click \textbf{Create a patron account}. Fill in the form.

After submitting, the screen shows a \textbf{TOTP secret} once and never again.

\proof{Your password is never stored. The database keeps a random 128-bit salt
and \texttt{SHA-256(salt + password)}. Your email, phone, card number, name and
address are each encrypted with RSA-2048 before being saved.}

\subsection*{F2. Browse the catalog}

Go to \textbf{Catalog}. Search for ``crypt''. Click a genre chip. Switch the
dropdown to ``Available now''.

\proof{Book data is \emph{not} secret --- the library must be searchable --- so
it is stored as normal text. Only \emph{your} activity is encrypted.}

\subsection*{F3. Borrow a book}

Open any book. Read the panel on the right, then press \textbf{Borrow this book}.

\proof{A record containing the book id, title, author and dates is encrypted
with ECIES (elliptic curve) and given an HMAC tag before it touches the
database.}

\subsection*{F6. My Vault}

Go to \textbf{My Vault}. You see the book, a due-date countdown, and a green
\textbf{Verified} badge.

Now press \textbf{Stored row} on that loan.

\proof{The app shows you the actual ciphertext and HMAC tag stored in the
database. The readable version above it exists only in memory --- it is never
saved in readable form.}

\subsection*{F4 + F5. Return a book and get fined}

To create an overdue book on demand, run this in a terminal:

\begin{verbatim}
   node scripts/backdate.mjs faiza 18
\end{verbatim}

Then go back to \textbf{My Vault} and press \textbf{Return}.

You should get a fine of exactly \textbf{185 BDT} plus a replacement notice:

\begin{center}
\begin{tabular}{lrrr}
\toprule
\textbf{Tier} & \textbf{Days} & \textbf{Rate} & \textbf{Subtotal} \\
\midrule
Days 1--7   & 7 & 5 BDT  & 35 BDT \\
Days 8--14  & 7 & 10 BDT & 70 BDT \\
Day 15+     & 4 & 20 BDT & 80 BDT \\
\midrule
\textbf{Total} & \textbf{18} & & \textbf{185 BDT} \\
\bottomrule
\end{tabular}
\end{center}

Now try to borrow another book --- you will be blocked. Press \textbf{Pay},
then try again --- it works.

\proof{The tiered fine rule from the proposal, the borrowing block while fines
are unpaid, and the replacement notice at 15+ days.}

\subsection*{F7. Private reviews}

Go to \textbf{Reviews} and write one. Only books you have returned appear in
the list.

\proof{The review text \emph{and the name of the book} are both inside the
ciphertext. The librarian cannot even tell which book you wrote about.}

\subsection*{F8. Secure chat}

Go to \textbf{Chat} and send a message to the librarian.

\proof{The message is encrypted with the librarian's RSA public key. A second
copy is encrypted with your own key so you can still read your own outbox. The
tag is an HMAC using a shared secret both sides compute with ECDH.}

\subsection*{F9. Profile}

Go to \textbf{Profile}. A green \textbf{CBC-MAC verified} banner sits at the
top. Press \textbf{Show} on the right panel to see the stored ciphertext.

Edit any field and save --- the app tells you which fields were re-encrypted.

\proof{Each field is separately RSA-encrypted, and one CBC-MAC tag covers the
whole profile so a field cannot be swapped or deleted without detection.}

\subsection*{F10. Key rotation}

Go to \textbf{Security}. Note the key version (\texttt{v1}) and the two
fingerprints. Press \textbf{Rotate my keys}, confirm with your password.

\proof{New RSA and ECC key pairs are made, every record you own is decrypted
with the old keys and re-encrypted with the new ones, and the old keys are
\emph{deleted}. The report shows the count per table and the version becomes
\texttt{v2}.}

% ==========================================================
\newpage
\section{Part B --- Librarian walkthrough}

Sign out, then sign in as \texttt{librarian}.

\subsection*{A5 + the main privacy proof: Overview}

The \textbf{Overview} page has a panel called \textbf{What the librarian can
see}. Use the four buttons to switch between checkouts, reviews, fines and
chat messages.

\proof{\textbf{This is the most important screen in the whole demo.} It is a
direct read of the database with full admin access --- and all you see is
hex. The librarian holds no key that opens it.}

\subsection*{A1. Book inventory}

Go to \textbf{Books}. Add a book, edit one, try to delete one that is on loan
(it will refuse).

\proof{Only the admin can write to the catalog. Patrons have read-only access.}

\subsection*{A2. Account management}

Go to \textbf{Patrons}. You can see usernames, join dates, last login and
record counts. Suspend a patron.

\proof{Notice what is \emph{missing}: no email, no phone, no reading history.
Suspension instantly kills that patron's sessions.}

\subsection*{A3. Fine dashboard}

Go to \textbf{Fines}. Click on a ciphertext cell in the last column.

\proof{The librarian sees the \emph{amount} (needed to run the library) but the
record body --- which book, which dates --- stays encrypted.}

\subsection*{A4. Reply to patrons}

Go to \textbf{Chat} and reply to the message you sent as the patron.

\subsection*{A5. Audit log}

Go to \textbf{Audit Log}. Set the outcome filter to \textbf{failure}, then
\textbf{warning}.

\proof{Logins, failed passwords, wrong OTP codes, RBAC refusals, tampering and
key rotations are all recorded --- but only metadata, never content.}

% ==========================================================
\section{Part C --- Showing the nine security features}

\renewcommand{\arraystretch}{1.35}
\begin{longtable}{@{}p{0.4cm} p{4.2cm} p{10.3cm}@{}}
\toprule
\textbf{\#} & \textbf{Feature} & \textbf{How to show it} \\
\midrule
\endfirsthead
\toprule
\textbf{\#} & \textbf{Feature} & \textbf{How to show it} \\
\midrule
\endhead

1 & Registration \& Login &
Register a new account. Then open Supabase $\rightarrow$ Table Editor
$\rightarrow$ \texttt{users}. The \texttt{password\_hash} column is a hash, and
\texttt{email\_enc} is unreadable hex. \\

2 & Two-Factor Authentication &
On the sign-in screen, enter the correct password and stop. You get a
challenge, \emph{not} a session. Enter a wrong code --- refused. Only the
correct code signs you in. \\

3 & Encrypted User Profile &
Profile page $\rightarrow$ \textbf{Show}. Compare the readable fields on the
left with the ciphertext on the right. \\

4 & Key Management Module &
Security page: fingerprints, key version, and the \textbf{Rotate my keys}
button with its report. \\

5 & Encrypted Posts &
Borrow a book and write a review, then look at the librarian's Overview panel
for those same rows. \\

6 & MAC Integrity Layer &
In My Vault, press \textbf{Simulate tampering}, then reload. The badge turns
red and decryption is refused. \emph{(See the box below.)} \\

7 & Role-Based Access Control &
While signed in as the patron, open
\texttt{localhost:4000/api/admin/users} in the same browser. You get a
\texttt{403 FORBIDDEN}, and an \texttt{rbac\_denied} entry appears in the audit
log. \emph{(See the box below --- this is the important version.)} \\

8 & Secure Chat &
Send a message, then check the librarian's Overview $\rightarrow$ Chat
messages. The body is hex. \\

9 & Secure Session Management &
Sign in as the patron in one browser, then sign in again in another. The first
one is dead --- one session per account. \\

\bottomrule
\end{longtable}

\subsection*{The tamper demo (feature 6) --- do this one slowly}

This is the clearest single demonstration in the project.

\begin{enumerate}[leftmargin=1.4cm, itemsep=3pt]
  \item In \textbf{My Vault}, note the loan shows a green \textbf{Verified} badge.
  \item Press \textbf{Simulate tampering}. This flips one character of the
        stored ciphertext in the database --- exactly what an attacker with
        database access would do.
  \item Reload the page.
  \item The row turns red: \textbf{Tampered --- MAC mismatch}. The book title
        is gone, and the \textbf{Return} button is disabled.
  \item Sign in as the librarian $\rightarrow$ \textbf{Audit Log}. An
        \texttt{integrity\_failure} warning has been recorded.
\end{enumerate}

\textbf{Why this matters:} the system checks the MAC \emph{before} attempting
to decrypt. So altered data is never turned into readable output --- it is
refused. That is called \emph{encrypt-then-MAC}.

\subsection*{The RBAC demo (feature 7) --- show the server, not the screen}

If you only type \texttt{localhost:5173/admin} as a patron, the page says
``not authorised'' --- but that is just the \emph{browser} hiding a menu. A
teacher may point out that hiding a page is not security.

So show the server refusing instead. While signed in as the patron, open this
address in the same browser:

\begin{verbatim}
   http://localhost:4000/api/admin/users
\end{verbatim}

You get back:

\begin{verbatim}
   {"error":{"code":"FORBIDDEN",
             "message":"This area is restricted to admin accounts."}}
\end{verbatim}

Now sign in as the librarian and open \textbf{Audit Log}. There is a new red
\texttt{rbac\_denied} row:

\begin{verbatim}
   GET /api/admin/users requires admin, caller is patron
\end{verbatim}

\textbf{Why this matters:} your session cookie was sent, you were properly
authenticated --- and the server still said no, because of your \emph{role}.
That is real access control, checked on the server every single request.

% ==========================================================
\newpage
\section{Part D --- Lab requirements: where each one lives}

\renewcommand{\arraystretch}{1.3}
\begin{longtable}{@{}p{0.4cm} p{6.2cm} p{8.3cm}@{}}
\toprule
\textbf{\#} & \textbf{Requirement} & \textbf{Where it is} \\
\midrule
\endfirsthead
\toprule
\textbf{\#} & \textbf{Requirement} & \textbf{Where it is} \\
\midrule
\endhead

1 & Login \& registration modules & \texttt{server/src/routes/auth.js} \\
2 & All user info encrypted before storage & RSA per field ---
    \texttt{server/src/services/records.js} \\
3 & Passwords hashed + salted & Custom SHA-256, 128-bit salt ---
    \texttt{server/src/crypto/password.js} \\
4 & Two-step authentication & Password, then HMAC-TOTP. Session issued only
    after both pass. \\
5 & Key Management Module & \texttt{server/src/services/keyManagement.js} \\
6 & Create / view / edit posts, auto-encrypted & Borrow, reviews, fines ---
    \texttt{routes/checkouts.js}, \texttt{reviews.js}, \texttt{fines.js} \\
7 & No plaintext in the database & \texttt{db/schema.sql} has no plaintext
    content column. Verify in Supabase. \\
8 & MAC for integrity & HMAC-SHA256 on records and chat, CBC-MAC on profile \\
9 & Only asymmetric encryption & RSA-2048 and ECIES over secp256k1. No AES
    anywhere. \\
10 & At least two asymmetric algorithms & RSA for profile and chat; ECC for
    records, reviews, fines \\
11 & RBAC & \texttt{server/src/middleware/auth.js}, enforced server-side \\
12 & Secure session management & \texttt{server/src/services/session.js} \\

\bottomrule
\end{longtable}

% ==========================================================
\section{Part E --- Proving the cryptography}

All algorithms are written by hand. Nothing uses Node's \texttt{crypto}
module or any crypto library.

\subsection*{1. Run the test suite}

\begin{verbatim}
   npm test
\end{verbatim}

\textbf{Result: 107 tests pass.} This checks:

\begin{itemize}[leftmargin=1.4cm, itemsep=2pt]
  \item SHA-256 against the official published values (FIPS 180-4)
  \item HMAC-SHA256 against the official RFC 4231 test cases
  \item The prime test against Carmichael numbers (numbers that fool weaker tests)
  \item RSA encrypt/decrypt, including the exact maximum message size
  \item Elliptic curve maths ($n \times G$ = point at infinity, and others)
  \item That tampered ciphertext is always rejected
  \item \textbf{A scan that fails the build if any file imports a crypto library}
\end{itemize}

That last one is the strongest single answer to ``did you really write it
yourself?''

\subsection*{2. Run the full system test}

\begin{verbatim}
   node scripts/smoke-test.mjs
\end{verbatim}

\textbf{Result: 79 checks pass.} It registers a user, signs in with both
factors, borrows, goes overdue, gets fined 185 BDT, pays, reviews, chats,
rotates keys, and confirms the database holds no readable data.

\subsection*{3. Live in the browser}

Open \texttt{http://localhost:4000/api/lab/self-test}

Runs the published test values live and returns \texttt{17/17}. No login
needed --- your teacher can open it directly.

% ==========================================================
\newpage
\section{Part F --- Likely questions and short answers}

\textbf{Q. Where is my password stored?} \\
Nowhere. Only a random salt and \texttt{SHA-256(salt + password)}. The password
itself is never written down.

\textbf{Q. Why two different asymmetric algorithms?} \\
RSA for small, occasional data (profile fields, chat messages). ECC/ECIES for
frequent, larger data (loans, reviews, fines) because it gives much smaller
ciphertext for the same strength.

\textbf{Q. Why no AES or any symmetric cipher?} \\
Requirement 9 forbids it. Everything is RSA or ECC.

\textbf{Q. Then how is CBC-MAC built without a block cipher?} \\
The 128-bit block function is \texttt{SHA-256(key + block)} truncated to 16
bytes, chained in CBC mode. A length prefix is added first, which blocks the
classic CBC-MAC forgery attack.

\textbf{Q. Can the librarian read patron data?} \\
No. Private keys are stored wrapped with a value derived from the patron's own
password. They are unwrapped only during that patron's session and kept in the
server's memory --- never written to disk.

\textbf{Q. What happens if the whole database leaks?} \\
The attacker gets ciphertext, MAC tags, salted hashes, and wrapped keys they
cannot unwrap without each user's password.

\textbf{Q. Why is \texttt{due\_date} readable but the book is not?} \\
The library needs due dates to work at all. But \texttt{book\_id} is
deliberately \emph{not} stored in readable form, so nobody with database access
can build a picture of what you read.

\textbf{Q. What stops someone replaying a chat message?} \\
The timestamp is included in the HMAC calculation. Change the time and the tag
no longer matches.

\textbf{Q. What is the session token?} \\
256 random bits. Only its SHA-256 hash is stored. The cookie is HttpOnly, tied
to your IP and browser, expires after 30 minutes idle, and only one session per
account is allowed.

\textbf{Q. Why is Row Level Security disabled in Supabase?} \\
Because the proposal says so. Supabase is used only as a database --- all
access control is written by hand in the Express backend, which is what
requirement 11 asks for. Turning on RLS would hand that job to Supabase.

\textbf{Q. Where does the randomness come from?} \\
\texttt{crypto.getRandomValues()} only --- the one operating-system call the
lab permits. It is used in exactly one file,
\texttt{server/src/crypto/random.js}, and the test suite enforces that.

\textbf{Q. What happens to old keys after rotation?} \\
They are deleted from the database, not archived. An archived key is a key that
can still be stolen.

\subsection*{Honest limitations (good to say before being asked)}

\begin{itemize}[leftmargin=1.4cm, itemsep=3pt]
  \item Encryption happens on the \textbf{server}, not in the browser. So the
        server sees data briefly while you are signed in. True end-to-end
        encryption would need the crypto to run in the browser.
  \item Restarting the server logs everybody out, because unwrapped keys live
        only in memory. That is the cost of never storing keys in usable form.
  \item The tamper demo button is a teaching tool. It is switched off by
        setting \texttt{ENABLE\_LAB\_ROUTES=false} in \texttt{server/.env}.
\end{itemize}

\vspace{1em}
\begin{center}
\fbox{\begin{minipage}{0.92\textwidth}
\centering
\textbf{One-line summary:} the librarian runs the library; the patron owns the
data. Every record is encrypted before storage and its integrity is checked
before it is ever read back.
\end{minipage}}
\end{center}

\end{document}
